\section{Les cours facultatifs de langues}\label{les-cours-facultatifs-de-langues}

\stanceinfo{\textbf{Congrès de la FCEG 2018} [2018-01-07]}{2024-01-02}

\officialstance{La Fédération canadienne étudiante de génie croit que les étudiants canadiens en génie devraient avoir l'occasion de suivre des cours facultatifs de langue dans le cadre de leur programme de premier cycle afin d'enrichir leur diplôme d'ingénierie.}

\stanceheading{La position des étudiants}
\begin{itemize}
 \item En 2018, le BCAPG a mis à jour ses \href{https://engineerscanada.ca/sites/default/files/accreditation/2021-2022-cycle/accreditation-criteria-procedures-2020.pdf}{critères et procédures d'accréditation} [2] afin d'accepter explicitement les cours facultatifs de langue dans le cadre de leurs exigences en matière « d'études complémentaires » pour les programmes d'ingénierie au Canada ; cependant, environ 23\% des écoles d'ingénieurs accréditées au Canada n'ont pas encore reflété cette évolution au sein de leurs facultés.
 \item En 2022, le critère a été supprimé 3.4.5.2 comme indiqué ci-dessous, où la FCEG a constaté qu'il reste inscrit sur certains sites d'universités d'études complémentaires :
 \begin{itemize}
  \item [Supprimé] Critère 3.4.5.2 L'enseignement des langues peut être inclus dans les études complémentaires à condition qu'il ne soit pas suivi pour remplir une condition d'admission. En outre, le contenu du programme d'études qui transmet principalement des compétences linguistiques peut être pris en compte dans l'unité d'accréditation (UA) requise pour les études complémentaires, mais ne peut pas être utilisé pour satisfaire aux exigences relatives aux matières qui traitent des questions centrales, des méthodologies et des processus de pensée des sciences humaines et sociales.
  \item Le FCEG est en accord avec la déclaration supprimée, car même les matières qui traitent d'autres langues sont inclusives pour l'enrichissement des étudiants en ingénierie en dehors de leurs cours techniques.
 \end{itemize}
 \item La connaissance d'une autre langue et des matières connexes est un atout précieux pour l'avenir d'un étudiant en ingénierie et les étudiants doivent disposer d'un éventail complet et approprié d'options linguistiques.
\end{itemize}

\stanceheading{Les recommandations aux partenaires, aux parties intéressées et aux autres entités}
\begin{itemize}
 \item La FCEG incite le BCAPG à aider les 23 \% d'établissements canadiens restants à interpréter leur politique actuelle sur les cours facultatifs en publiant une déclaration interprétative officielle sur le sujet, énonçant des critères supplémentaires.
 \item La FCEG incite les DDIC à encourager ses facultés membres à revoir et à reconsidérer leurs offres de cours facultatifs, et à permettre aux étudiants en ingénierie de s'inscrire à des cours de langue qui comptent pour l'obtention de leur diplôme, en plus et y compris des deux langues officielles du Canada.
 \item La FCEG incite les DDIC à mettre à jour ses politiques et ses informations disponibles sur ses calendriers académiques afin de suivre et de décrire avec précision les politiques et procédures du BCAPG concernant ce qui est considéré comme un cours facultatif.
 \item La FCEG incite les associations des étudiants en génie à soutenir les efforts de plaidoyer nationaux et à tirer parti de la mise à jour de la politique du BCAPG pour défendre la mise en œuvre de cours facultatifs de langue dans leur programme afin de profiter des avantages du multilinguisme dans un pays diversifié et un monde de plus en plus globalisé.
 \item La FCEG incite ses membres à remettre en question le cursus actuellement en place, et à poursuivre des études qui les inspirent en langues et qui enrichissent leur formation d'ingénieur, même si des autorisations supplémentaires sont nécessaires pour les poursuivre.
\end{itemize}

\stanceheading{Ce que fait la FCEG}
\begin{itemize}
 \item Les discussions de la FCEG sur l'interprétation de la politique des cours facultatifs de langues sur deux périodes de mandat ont été acceptées par le Comité des politiques et procédures du BCAPG qui a accepté que les sections linguistiques soient considérées comme des cours à option ouverts avant la Réunion du Président (RdP) de 2019.
 \item La FCEG a créé et continue de mettre à jour un dossier de plaidoyer pour aider les étudiants à plaider en faveur des cours facultatifs de langues au niveau de l'établissement.
 \item La FCEG a examiné les options linguistiques au sein des universités d'ingénierie accréditées au Canada et a dressé une liste des universités qui ne respectent pas la politique actuelle.
 \item La FCEG a présenté des données sur les universités BCAPG qui ne respectent pas la politique actuelle ou qui utilisent des informations obsolètes afin de faciliter le dialogue entre les DDIC et le BCAPG.
\end{itemize}

\stanceheading{Ce que la FCEG envisage de faire}
\begin{itemize}
 \item La FCEG fournira des ressources sur les cours facultatifs de langues afin de fournir à ses écoles membres les informations dont elles ont besoin pour plaidoyer en faveur de la modification des cours facultatifs dans leurs établissements respectifs.
 \item La FCEG fournira une analyse des avantages des cours facultatifs de langues disponibles comme études complémentaires et de la manière dont ils enrichissent la formation en ingénierie.
 \item La FCEG travaillera avec les DDIC pour lui faire connaître les politiques et les documents qui ne correspondent pas à la politique et aux procédures actuelles du BCAPG ou qui ne sont pas tout à fait exacts.
 \item La FCEG poursuivra ses efforts de plaidoyer et continuera la discussion sur l'importance du bilinguisme/multilinguisme lors des présentations des partenaires externes sur l'avenir de la formation des ingénieurs.
\end{itemize}

\newpage
